\documentclass[11pt,a4paper]{article}

\usepackage{amsmath,amssymb,amsthm}
\usepackage{graphicx}
\usepackage{hyperref}
\usepackage{algorithm}
\usepackage{algpseudocode}
\usepackage{booktabs}
\usepackage{geometry}
\usepackage{authblk}

\geometry{margin=1in}

\theoremstyle{definition}
\newtheorem{definition}{Definition}
\newtheorem{theorem}{Theorem}
\newtheorem{lemma}{Lemma}
\newtheorem{proposition}{Proposition}
\newtheorem{corollary}{Corollary}

\title{Algorithmic Stablecoins on Lux: \\
Stability Mechanisms and Peg Maintenance}

\author{Zach Kelling}
\affil{Hanzo Industries \\ Lux Industries \\ Zoo Labs Foundation}
\date{2022}

\begin{document}

\maketitle

\begin{abstract}
We present a rigorous analysis of algorithmic stablecoin mechanisms deployable on the Lux blockchain. Our framework covers three primary stability approaches: over-collateralized debt positions, algorithmic supply adjustment, and hybrid fractional-reserve systems. We derive stability conditions, analyze failure modes, and provide formal proofs of peg maintenance under various market conditions. The analysis demonstrates that robust stablecoin design requires careful consideration of collateral quality, liquidation efficiency, and incentive alignment. We conclude with implementation guidelines optimized for Lux's high-throughput environment.
\end{abstract}

\section{Introduction}

Stablecoins are digital assets designed to maintain a stable value relative to a reference asset, typically the US dollar. They serve as the primary medium of exchange in decentralized finance, enabling users to hold value without exposure to cryptocurrency volatility.

Three primary mechanisms exist for maintaining price stability:

\begin{enumerate}
\item \textbf{Custodial}: Backed 1:1 by fiat reserves (centralized)
\item \textbf{Over-Collateralized}: Backed by cryptocurrency collateral exceeding debt value
\item \textbf{Algorithmic}: Supply adjusted via economic incentives
\end{enumerate}

This paper focuses on decentralized approaches (2) and (3), providing mathematical foundations for their design and analysis.

\section{Over-Collateralized Stablecoins}

\subsection{Collateralized Debt Positions}

\begin{definition}[Collateralized Debt Position (CDP)]
A CDP is a smart contract allowing users to:
\begin{enumerate}
\item Deposit collateral $C$
\item Mint stablecoin debt $D$
\item Subject to collateralization ratio $\rho = C/D > \rho_{\min}$
\end{enumerate}
\end{definition}

\begin{definition}[Collateralization Ratio]
For collateral value $C$ and debt $D$ in reference currency:
\begin{equation}
\rho = \frac{C}{D}
\end{equation}
The minimum collateralization ratio $\rho_{\min} > 1$ ensures over-collateralization.
\end{definition}

\subsection{Stability Analysis}

\begin{theorem}[Over-Collateralization Stability]
An over-collateralized stablecoin maintains its peg if:
\begin{equation}
\sum_i C_i > \sum_i D_i
\end{equation}
where $C_i$ and $D_i$ are collateral and debt for position $i$.
\end{theorem}

\begin{proof}
If total collateral exceeds total debt, rational arbitrageurs can:
\begin{itemize}
\item If price $< \$1$: Buy stablecoin, redeem for $>\$1$ collateral
\item If price $> \$1$: Mint new stablecoin (if $\rho > \rho_{\min}$), sell for $>\$1$
\end{itemize}
Arbitrage drives price toward $\$1$.
\end{proof}

\subsection{Liquidation Mechanism}

\begin{definition}[Liquidation Trigger]
A CDP is liquidatable when:
\begin{equation}
\rho = \frac{C}{D} < \rho_{\text{liq}}
\end{equation}
where $\rho_{\text{liq}}$ is the liquidation threshold.
\end{definition}

\begin{proposition}[Liquidation Buffer]
The buffer between minimum collateralization and liquidation:
\begin{equation}
\text{Buffer} = \rho_{\min} - \rho_{\text{liq}}
\end{equation}
provides time for users to add collateral before liquidation.
\end{proposition}

\subsection{Stability Fee}

\begin{definition}[Stability Fee]
The stability fee $r_s$ is interest charged on outstanding debt:
\begin{equation}
D(t) = D_0 \cdot e^{r_s \cdot t}
\end{equation}
\end{definition}

\begin{theorem}[Stability Fee Impact]
The stability fee affects stablecoin supply:
\begin{itemize}
\item $r_s \uparrow \implies$ Borrowing cost increases $\implies$ Supply decreases
\item $r_s \downarrow \implies$ Borrowing cost decreases $\implies$ Supply increases
\end{itemize}
\end{theorem}

This provides a monetary policy lever for peg maintenance.

\section{Algorithmic Stablecoins}

\subsection{Seigniorage Model}

\begin{definition}[Two-Token Seigniorage]
A seigniorage stablecoin system consists of:
\begin{itemize}
\item Stablecoin $S$ with target price $P^* = \$1$
\item Governance/equity token $G$ absorbing volatility
\end{itemize}
\end{definition}

\subsection{Mint-Burn Mechanism}

\begin{algorithm}
\caption{Seigniorage Mint/Burn}
\begin{algorithmic}[1]
\If{$P_S > P^* + \epsilon$} \Comment{Above peg}
\State User deposits $\$1$ worth of $G$
\State Protocol mints 1 $S$ to user
\State $G$ is burned
\ElsIf{$P_S < P^* - \epsilon$} \Comment{Below peg}
\State User deposits 1 $S$
\State Protocol mints $\$1$ worth of $G$ to user
\State $S$ is burned
\EndIf
\end{algorithmic}
\end{algorithm}

\subsection{Stability Analysis}

\begin{theorem}[Seigniorage Stability Condition]
The seigniorage mechanism maintains peg if:
\begin{equation}
\text{Market Cap}(G) > \text{Supply}(S) \cdot |\Delta P_S|_{\max}
\end{equation}
where $|\Delta P_S|_{\max}$ is the maximum expected price deviation.
\end{theorem}

\begin{proof}
When $S$ trades below peg, users burn $S$ for $G$. This requires:
\begin{enumerate}
\item Users trust $G$ maintains value
\item Sufficient $G$ liquidity to absorb redemptions
\end{enumerate}
If $\text{MC}(G) < \text{Supply}(S)$, a bank run depletes $G$ value before restoring peg.
\end{proof}

\subsection{Death Spiral Dynamics}

\begin{definition}[Death Spiral]
A reflexive feedback loop where:
\begin{enumerate}
\item $P_S < P^* \implies$ Users redeem $S$ for $G$
\item $G$ supply increases $\implies P_G$ decreases
\item Confidence in system falls $\implies P_S$ decreases further
\item Repeat until system collapse
\end{enumerate}
\end{definition}

\begin{theorem}[Death Spiral Prevention]
A seigniorage system avoids death spirals if:
\begin{equation}
\frac{d(\text{MC}(G))}{d(\text{Supply}(S))} > -1
\end{equation}
This ensures redemptions don't destroy more value than they remove.
\end{theorem}

\section{Fractional-Reserve Hybrid}

\subsection{Design Principles}

\begin{definition}[Fractional Stablecoin]
A hybrid system combining:
\begin{itemize}
\item Fraction $\alpha$ backed by exogenous collateral
\item Fraction $1 - \alpha$ backed by governance token
\end{itemize}
\end{definition}

\begin{equation}
1 \text{ Stablecoin} = \alpha \cdot \$1_{\text{collateral}} + (1 - \alpha) \cdot \$1_{\text{governance}}
\end{equation}

\subsection{Dynamic Collateral Ratio}

\begin{definition}[Adaptive Collateral Ratio]
The collateral ratio $\alpha$ adjusts based on peg stability:
\begin{equation}
\frac{d\alpha}{dt} = k \cdot (P^* - P_S)
\end{equation}
where $k > 0$ is adjustment speed.
\end{definition}

\begin{itemize}
\item $P_S < P^*$: Increase $\alpha$ (more collateral backing)
\item $P_S > P^*$: Decrease $\alpha$ (more algorithmic)
\end{itemize}

\subsection{Stability Properties}

\begin{theorem}[Fractional Stability]
A fractional stablecoin with $\alpha \in (0, 1)$ is more stable than pure algorithmic ($\alpha = 0$) when:
\begin{equation}
\text{Volatility}(G) > \text{Volatility}(\text{Collateral})
\end{equation}
\end{theorem}

\begin{proof}
Redemption value:
\begin{equation}
V = \alpha \cdot P_{\text{collateral}} + (1 - \alpha) \cdot P_G
\end{equation}
Variance:
\begin{equation}
\text{Var}(V) = \alpha^2 \text{Var}(P_C) + (1-\alpha)^2 \text{Var}(P_G) + 2\alpha(1-\alpha)\text{Cov}(P_C, P_G)
\end{equation}
For $\text{Var}(P_G) > \text{Var}(P_C)$ and low correlation, increasing $\alpha$ reduces variance.
\end{proof}

\section{Peg Stability Mechanisms}

\subsection{Peg Stability Module (PSM)}

\begin{definition}[Peg Stability Module]
A PSM enables direct 1:1 swaps between stablecoin and approved collateral (e.g., USDC):
\begin{align}
\text{Swap In}: & \quad 1 \text{ USDC} \rightarrow 1 \text{ Stablecoin} - \text{fee} \\
\text{Swap Out}: & \quad 1 \text{ Stablecoin} \rightarrow 1 \text{ USDC} - \text{fee}
\end{align}
\end{definition}

\begin{theorem}[PSM Peg Bound]
With PSM fee $f$, the stablecoin price is bounded:
\begin{equation}
P_{\text{USDC}} \cdot (1 - f) \leq P_S \leq P_{\text{USDC}} \cdot (1 + f)
\end{equation}
\end{theorem}

\begin{proof}
Arbitrage opportunities:
\begin{itemize}
\item $P_S > P_{\text{USDC}}(1 + f)$: Deposit USDC, mint stablecoin, sell
\item $P_S < P_{\text{USDC}}(1 - f)$: Buy stablecoin, redeem for USDC
\end{itemize}
\end{proof}

\subsection{Interest Rate Control}

\begin{definition}[Savings Rate]
The stablecoin savings rate $r_{\text{save}}$ incentivizes holding:
\begin{equation}
\text{Staked Balance}(t) = \text{Balance}_0 \cdot e^{r_{\text{save}} \cdot t}
\end{equation}
\end{definition}

\begin{proposition}[Rate Policy]
\begin{itemize}
\item $P_S < P^*$: Increase $r_{\text{save}}$ to increase demand
\item $P_S > P^*$: Decrease $r_{\text{save}}$ to reduce demand
\end{itemize}
\end{proposition}

\section{Oracle Design}

\subsection{Price Feed Requirements}

Stablecoins require accurate price feeds for:
\begin{enumerate}
\item Collateral valuation
\item Liquidation triggers
\item Peg deviation detection
\end{enumerate}

\begin{definition}[Oracle Security Model]
A secure oracle provides price $\hat{P}$ such that:
\begin{equation}
|P - \hat{P}| < \epsilon
\end{equation}
with probability $1 - \delta$ for acceptable deviation $\epsilon$ and failure rate $\delta$.
\end{definition}

\subsection{Oracle Manipulation Resistance}

\begin{theorem}[Manipulation Cost]
For a TWAP oracle with window $w$ and underlying liquidity $L$:
\begin{equation}
\text{Cost to manipulate by } \Delta P \geq \Delta P \cdot L \cdot w
\end{equation}
\end{theorem}

\section{Implementation on Lux}

\subsection{Architecture}

\begin{enumerate}
\item \textbf{Core Module}: CDP management, minting, burning
\item \textbf{Liquidation Engine}: Monitors and executes liquidations
\item \textbf{Oracle Module}: Aggregates price feeds
\item \textbf{PSM Module}: Direct stablecoin swaps
\item \textbf{Governance}: Parameter updates via token voting
\end{enumerate}

\subsection{CDP Creation}

\begin{algorithm}
\caption{Open CDP}
\begin{algorithmic}[1]
\Require User $u$, collateral amount $C$, mint amount $D$
\Ensure CDP created or reverted
\State $P_C \gets \text{oracle.getPrice}(\text{collateral})$
\State $\rho \gets \frac{C \cdot P_C}{D}$
\State \textbf{require} $\rho \geq \rho_{\min}$
\State \textbf{require} $D + \text{totalDebt} \leq \text{debtCeiling}$
\State \textbf{transfer}(\text{collateral}, $u$, \text{vault}, $C$)
\State \textbf{mint}(\text{stablecoin}, $u$, $D$)
\State \textbf{emit} CDPOpened($u$, $C$, $D$)
\end{algorithmic}
\end{algorithm}

\subsection{Gas Optimization}

\begin{itemize}
\item \textbf{Batched Liquidations}: Multiple CDPs in single transaction
\item \textbf{Lazy Interest}: Compute on interaction, not per-block
\item \textbf{Proxy Pattern}: Upgradeable contracts with minimal storage
\end{itemize}

\section{Risk Analysis}

\subsection{Black Swan Events}

\begin{definition}[Black Swan Risk]
Scenarios where collateral value drops faster than liquidation can process:
\begin{equation}
\frac{dP_C}{dt} < -\frac{\rho_{\text{liq}} - 1}{\Delta t_{\text{liq}}}
\end{equation}
where $\Delta t_{\text{liq}}$ is liquidation processing time.
\end{definition}

\subsection{Mitigation Strategies}

\begin{enumerate}
\item \textbf{Emergency Shutdown}: Freeze system, allow collateral claims
\item \textbf{Debt Auctions}: Mint governance tokens to cover shortfall
\item \textbf{Insurance Fund}: Reserve buffer from stability fees
\end{enumerate}

\section{Conclusion}

We have presented a comprehensive framework for algorithmic stablecoin design on Lux:

\begin{enumerate}
\item \textbf{Over-Collateralized}: Mathematically guaranteed backing
\item \textbf{Seigniorage}: Efficient but requires confidence
\item \textbf{Hybrid}: Balanced risk-efficiency tradeoff
\end{enumerate}

Successful stablecoin design requires:
\begin{itemize}
\item Conservative collateralization ratios
\item Efficient liquidation mechanisms
\item Reliable oracle infrastructure
\item Robust governance for parameter updates
\end{itemize}

Lux's high throughput enables responsive liquidations and real-time peg maintenance, making it an ideal platform for stablecoin deployment.

\begin{thebibliography}{9}
\bibitem{maker}
MakerDAO. (2017). The Dai Stablecoin System. MakerDAO Whitepaper.

\bibitem{terra}
Kereiakes, E., Do, K., Di, D., Platias, N. (2019). Terra Money: Stability and Adoption. Terra Whitepaper.

\bibitem{frax}
Kazemian, S., et al. (2020). Frax: Fractional-Algorithmic Stablecoin Protocol. Frax Whitepaper.

\bibitem{liquity}
Kolten, R. (2021). Liquity: Decentralized Borrowing Protocol. Liquity Whitepaper.

\bibitem{rai}
Reflexer Labs. (2021). RAI: A Low Volatility, Trust Minimized Collateral. Reflexer Whitepaper.
\end{thebibliography}

\end{document}
